\section{课后习题：电像法}

% 循迹之外独立编号
\setcounter{example}{0}
\renewcommand{\theexample}{\arabic{example}}

\begin{example}[\itr{Surface Charge Distribution on a Grounded Conducting Sphere}{接地导体球的表面电荷分布} ---\refleaftext{solution10.1}]
A \itr{grounded conducting sphere}{接地导体球} of radius $R$ has a point charge $q$ located at a distance $d$ from its center ($d > R$). 
Find the \itr{surface charge density}{电荷面密度} $\sigma_e(\theta)$ as a function of the \itr{polar angle}{极角} $\theta$. 
\end{example}

\begin{example}[\itr{Potential outside a Grounded Infinite Cylindrical Conductor}{接地无限长圆柱形导体的电势分布} ---\refleaftext{solution10.2}]
An infinite \itr{grounded cylindrical conductor}{接地圆柱形导体} of radius $R$ is placed parallel to an infinite \itr{charged line}{带电直线} with \itr{linear charge density}{线电荷密度} $\lambda$ at a distance $a$ ($a > R$) from the cylinder axis. Use the \itr{method of images}{电像法} to find the \itr{electric potential}{电势} in the region outside the cylinder.
\end{example}

\begin{example}[\itr{Image of Image Charge Method}{像电荷的像电荷法} ---\refleaftext{solution10.3}]
An infinitely large horizontal \itr{grounding plane}{接地平面} at $z=0$ separates the entire space into two parts. In the upper space ($z>0$), there is a \itr{metallic sphere}{金属导体球} of radius $R$ with total charge $Q$ centered at $(0,0,d)$ ($d > R$).

\begin{singlefigure}[]{electroimage_11.png}[0.70]
\end{singlefigure}

\begin{enumerate}[label=(\alph*)]
    \item How many \itr{boundary conditions}{边界条件} are needed? Write down the \textbf{equations} for these boundary conditions.

    \item Assuming the charge on the sphere is \textbf{initially uniformly distributed} on the surface, the straightforward way to put image charge would be another uniformly charged spherical shell of radius $R$ with total charge $-Q$ in the lower space. However, we can also use \itr{point image charges}{点像电荷} instead of a uniformly charged sphere image to obtain the \textbf{same results}. Find out the \textbf{minimum number} of image point charges required, and their \textbf{positions} $\vec{r}_i$ and \textbf{charges} $q_i$.

    \item Once we have the image charges in (b), the charge must \itr{redistribute}{重新分布} on the surface of the sphere. To conform with the requirement by the sphere, we can assume another image charge inside it. Calculate the \textbf{position} and \textbf{magnitude} of this ``\itr{image charge of image charge}{像电荷的像电荷}''. Then, in order to conform with the requirement by the plane, there will be another image charge below the plane.

    \item How many image charges are needed to \textbf{exactly} reproduce the boundary conditions? If possible, write down a \textbf{recursive expression} for the positions and magnitudes of these image charges.
\end{enumerate}
\end{example}
